\section{Method Development Results}
\label{sec:development_results}

All four planner arms completed all 60 development evaluations, with no
sample above 15~N. Table~\ref{tab:component_ablation_dev} shows that adaptive
objective and compute each recover part of the fixed-planner gap, while their
combination gives the lowest mean relative error (MRE) and normalised
root-mean-square error (NRMSE). The complete method reduces 12~N RMSE from
2.927 to 2.600~N and raises mean force from 9.308 to 9.714~N, although the
strict 12~N tracking gate remains 0/5 for every arm.

\begin{table*}[t]
\centering
\caption{Matched planner-component development ablation. MRE and NRMSE are
averaged over 15 seed--target evaluations per arm.}
\label{tab:component_ablation_dev}
\small
\begin{tabular}{lrrrr}
\toprule
Arm & MRE (\%) & NRMSE (\%) & Samples & Median time (ms) \\
\midrule
M1: fixed objective, fixed compute & 12.03 & 14.66 & 128.0 & 112.18 \\
M2: adaptive objective             & 10.94 & 13.77 & 128.0 & 107.86 \\
M3: adaptive compute               & 10.26 & 13.52 & 121.2 &  80.58 \\
M4: complete method                &  9.65 & 12.94 & 120.4 &  83.25 \\
\bottomrule
\end{tabular}
\end{table*}

A behaviour-cloning (BC) coefficient of 2 completes all 15 development runs
without a sampled violation; removing BC prevents contact in all 15, whereas
coefficient 0.5 yields 12/15 task passes and two samples above 15~N. Router
thresholds change planner use smoothly rather than revealing a discrete
operating point. The stable MPPI segment contains most of the 12~N negative
bias, and alternate risk representations do not improve the frozen
force-only representation. Complete development tables and sample-level
decompositions are provided in the Supplementary Material.

